%% ============================================================
\section{Worked Example: Null Dereference Detection}
\label{sec:worked-example}

To make the analysis pipeline concrete, consider this C function:

\begin{lstlisting}
void process(FILE *log) {
    char *buf = malloc(256);
    if (buf == NULL) {
        fprintf(log, "alloc failed\n");
        return;
    }
    /* buf is NotNull here */
    strcpy(buf, "initialized");
    free(buf);
}
\end{lstlisting}

SqC's analysis proceeds through the four stages:

\textbf{Stage 1 (Parse):} Tree-sitter produces a CST.  The
\texttt{malloc} call, \texttt{if} statement, \texttt{return},
\texttt{strcpy}, and \texttt{free} are all represented as nodes.

\textbf{Stage 2 (Prescan):} If cross-file analysis is enabled,
\texttt{malloc} is recognized from the standard function database
as a function that may return NULL.

\textbf{Stage 3 (CFG + Dataflow):} The CFG has four basic blocks:
entry (malloc call), the null-check branch, the error path (fprintf +
return), and the success path (strcpy + free).  The null-state dataflow
computes:
\begin{itemize}
    \item After \texttt{malloc}: \texttt{buf} $=$ PossiblyNull
    \item True branch (\texttt{buf == NULL}): \texttt{buf} $=$ DefinitelyNull
    \item False branch: \texttt{buf} $=$ NotNull
    \item At \texttt{strcpy}: \texttt{buf} $=$ NotNull (the DefinitelyNull
        path returned early)
\end{itemize}

\textbf{Stage 4 (Rule Evaluation):} EXP34-C checks every dereference
of a pointer.  At the \texttt{strcpy} call, \texttt{buf} is NotNull,
so no violation is reported.  If the \texttt{if (buf == NULL)} check
were removed, \texttt{buf} would still be PossiblyNull at \texttt{strcpy},
and EXP34-C would flag the violation.
